% !TEX root = ../main.tex
\chapter{Дүгнэлт}\label{ch:concl}

“Multithreaded Sorting Application” төслийг Java 17 дээр бүрэн хэрэгжүүлэв. Програм жагсаалтыг хоёр хуваан хоёр
sorting thread-ээр параметрээр өгсөн мужийг эрэмбэлж, merging thread-ээр хоёр дахь global массив руу нэгтгэнэ. Parent
thread бүх thread-ийг \code{join()} хийсний дараа үр дүнг хэвлэнэ. Үүн дээр таван алгоритм, оролтын бүрэн шалгалт,
JSON event trace, бодит thread-тэй холбогдсон Sorting animation (Replay ба Lockstep горим) хийж, 100 автомат тестээр
шалгасан.

Ажлын гол сургамж нь гурав. Нэгдүгээрт, \code{start()}/\code{join()} нь өөрөө синхрончлол бөгөөд давхцахгүй муж ба
join-ийн дарааллын ачаар lock огт хэрэггүй. Хоёрдугаарт, зэрэгцээ нь үнэгүй биш: жижиг $n$-д thread-ийн зардал
давамгайлж, их $n$-д ч нэгтгэх дан үеэс болж хурдсалт 2-т хүрэхгүй (Амдалын хууль). Гуравдугаарт, хурдсалтыг зөв
тайлбарлахын тулд “1 thread, хувааж” гэх мэт тохирох суурь хэмжилт болон олон JVM-ийн медиан шаардлагатай.
Цаашид олон thread бүхий мод хэлбэрийн нэгтгэл, зэрэгцээ merge ашиглаж олон цөмийг илүү бүрэн ашиглах боломжтой.
